\chapter{CONCLUSION}

The objectives of recognizing a hand-drawn circuit diagram, transforming it into an editable schematic having simulation capabilities, and analyzing said circuit with LLM-powered Q\&A have been successfully achieved. Functional pipelines for each aspect of the project have been designed and implemented effectively. For component detection, YOLOv7 has achieved a mAP@0.5 of 95\% and for text recognition, TrOCR has achieved accuracy of 84.5\%. Through the use of junction-based geometric inference and collinear merging, the wire detection algorithm extracts circuit connectivity. Overall, all these elements work hand-in-hand to extract components with their values and connections, effectively digitizing a circuit. Hence, the project has successfully addressed a significant gap in educational circuit analysis by making circuit simulation easier than ever.

Beyond the ML inference side of things, the project also pulls off a working full-stack deployment: a React.js frontend talks to a FastAPI backend, which in turn hooks into CircuitJS for running simulations.

CIRCUITRON also doubled as a research exercise for the authors. Multiple approaches for solving the same problem were tried out (different YOLO models for component detection, custom OCR and TrOCR for text recognition, and junction-based and path finding algorithms for wire connectivity) and the better alternative was selected in each case with limitations of the discarded method explained.
 
